\section{Authoring-Oriented Design Space}
\label{sec:design-space}

Existing taxonomies characterize recurring structures and relationships in
completed composite visualizations~\cite{javed2012exploring,Composite,vaid}.
To make these structures authorable, we operationalize them around three
decisions: what visual units should be composed, what data and spatial
relationships define their composition, and how these relationships are
realized across coordinate systems. We ground the composition units in expert
perception and use established composition patterns to derive their
relationships and spatial realizations, forming the design basis of {\system}.


\subsection{Perceptually Grounded Chart Blocks}

To derive units that remain visually recognizable yet independently
composable, we conducted an expert-grounded analysis of the VAID
corpus~\cite{vaid}, which contains 442 composite visualizations. We selected
the 100 visualizations with the longest specifications, using specification
length as a proxy for relatively complex composition structures. Six
visualization experts (E1--E6) participated, including four academic
researchers (E1--E4) and two industry practitioners (E5--E6). All had more
than five years of visualization experience and had designed or implemented
composite visualizations. E1 and E3 additionally had prior experience with
visualization surveys and classification.

\textbf{Block granularity.}
The experts decomposed the selected visualizations into units they considered
visually coherent and meaningful to construct or manipulate independently, and
we asked why these units should not be defined at a coarser or finer level.
Their responses suggested that block granularity is shaped jointly by
structural independence, established chart semantics, and perceptual identity.
Experts tended to preserve recognizable chart forms as single units even when
they could be decomposed into repeated graphical elements; for example, a
stacked bar chart was viewed as one chart rather than a composition of multiple
bar charts, and a sunburst as one chart rather than a composition of arcs.
At the same time, visualization tools and design practice have established
named idioms such as grouped and stacked bar charts, or streamgraphs and
stacked area charts, whose distinctions are already meaningful to authors.
We therefore define blocks around perceptually complete and conventionally
recognized chart structures, avoiding both coarse composite-level templates
and fine-grained decomposition into marks.




\textbf{Block families.}
To organize chart blocks for composite authoring, we examined what cues experts
use to identify the chart types they are looking for and to recognize related
design alternatives. The experts then performed a grouping task, organizing
blocks they expected to find under the same category in an authoring tool or
considered plausible alternatives. We converted their co-grouping judgments
into pairwise similarities and combined hierarchical clustering, majority
agreement, and qualitative rationales to identify recurring perceptual
relationships.

Three factors repeatedly shaped these judgments. Coordinate-system alternatives
can be structurally related yet perceptually distinct: most experts treated
icicle and sunburst charts as separate idioms, while E1 associated them as
alternatives. Dimensional alternatives, such as basic,
grouped, and stacked bar charts, were commonly related despite different data
requirements and mark organizations. Recognizable visual variants also
remained independently identifiable: E1 and E3 viewed step-line charts mainly
as line-chart variants, while E6 emphasized their distinct identity in some
domain-specific uses. We use these recurring perceptual cues to guide the
organization of block families, connecting related blocks while preserving
their explicit visual identities.



\subsection{Dimension-Aware Composition Semantics}

Existing taxonomies define composition primarily through the resulting visual
relationships~\cite{javed2012exploring,Composite,vaid}. For authoring, we
further characterize each composition type by how it organizes data dimensions
across chart blocks. This representation allows composition compatibility to
be checked from block specifications and connects chart-level data assignments
with composition-level structure.

\textbf{Layering} combines blocks within a shared coordinate space. The
participating blocks retain their own marks and encodings while sharing
compatible positional dimensions and scales. Layering therefore combines
visual structures without introducing an additional data dimension for layout.
\textbf{Concatenation} preserves separate chart blocks while establishing a
shared spatial reference along one dimension. The blocks may encode different
measures or structures, but their compatible positional dimension provides the
basis for alignment and coordinated placement.
\textbf{Faceting} promotes a data dimension to the composition level. Rather
than encoding that dimension within a single chart, the system partitions the
data and instantiates the same block for each value, using the promoted
dimension to organize the repeated views spatially.
\textbf{Nesting} similarly externalizes a data relationship into composition,
but associates child blocks with repeated elements of a parent block. A shared
key, filtering condition, or hierarchical correspondence determines which
child instance belongs to each parent element and provides its spatial anchor.

From this perspective, composition is not only a spatial arrangement of
charts, but also a way of reorganizing how data dimensions are represented.
In particular, dimensions can remain within a block, be promoted to repeated
views through faceting, or be expressed through parent--child relationships
in nesting. These semantics provide the basis for {\system}'s data-aware
composition and dimensional adaptation.


\subsection{Composition through Spatial References}

Composition is not determined by a block's coordinate system alone. Instead,
{\system} considers the spatial references exposed by a block and whether they
can support the relationship required by a composition. These references may
be explicit coordinate dimensions, derived structural axes, or positions
associated with repeated elements.

\textbf{Coordinate dimensions.}
Cartesian and polar blocks expose explicit spatial dimensions. Cartesian
blocks provide $x$ and $y$, while polar blocks provide angular ($\theta$) and
radial ($r$) dimensions. Layering can share both dimensions, whereas
concatenation shares one dimension and allocates the other. In polar
coordinates, this produces angular concatenation, which partitions $\theta$
while sharing $r$, and radial concatenation, which partitions $r$ while
sharing $\theta$. Geographic blocks similarly share a common geographic
reference, allowing compatible layers to be aligned through the same map
projection.

\textbf{Other spatial references.}
Useful spatial references may also arise from visual structure:
1) Repeated or ordered elements can expose derived
alignment axes; for example, the ordered leaves of a tree can form an axis even
when the tree layout has no conventional data axis, while nodes arranged around
a chord diagram provide angular positions that can align with repeated views.
2) Individual marks or structural elements can instead serve as element-level
anchors. Nesting uses such anchors to position child blocks and derive their
data context, allowing parent and child blocks to use different coordinate
systems. Thus, layouts such as force-directed graphs, bubble layouts, Venn
diagrams, and word clouds may lack semantically meaningful axes while still
providing positions that support nesting. 3) Finally, some blocks derive their
geometry from another block rather than defining an independent spatial
reference. Network links, for example, obtain their endpoints from node
positions and therefore inherit the spatial reference of the node layer,
allowing them to participate in layering wherever the referenced nodes are
positioned.

